% BHU PYQ -- Manually transcribed & error-corrected
\documentclass[11pt,a4paper]{article}

\usepackage[a4paper,top=2.5cm,bottom=2.5cm,left=2.8cm,right=2.8cm,headheight=14pt]{geometry}
\usepackage{amsmath,amssymb}
\usepackage[T1]{fontenc}
\usepackage[utf8]{inputenc}
\usepackage{lmodern}
\usepackage{microtype}
\usepackage{fancyhdr}
\usepackage{enumitem}
\usepackage{hyperref}
\usepackage{lastpage}
\usepackage{parskip}

\hypersetup{colorlinks=true,urlcolor=black,linkcolor=black,
  pdftitle={BPT-201: Thermal Physics -- BHU PYQ},pdfauthor={Banaras Hindu University}}

\pagestyle{fancy}\fancyhf{}
\renewcommand{\headrulewidth}{0.4pt}\renewcommand{\footrulewidth}{0.4pt}
\fancyhead[L]{\small\textit{Banaras Hindu University}}
\fancyhead[R]{\small\textit{BPT-201: Thermal Physics}}
\fancyfoot[C]{\small Page \thepage\ of \pageref{LastPage}}

\newlist{parts}{enumerate}{2}
\setlist[parts,1]{label=(\alph*),leftmargin=*,itemsep=5pt,topsep=3pt}
\setlist[parts,2]{label=(\roman*),leftmargin=*,itemsep=3pt,topsep=2pt}
\newcommand{\pts}[1]{\hfill\textit{[#1]}}

\begin{document}

\begin{center}
  {\large\bfseries BANARAS HINDU UNIVERSITY}\\[2pt]
  {\normalsize B.Sc. (Hons.) Semester II Examination, 2013-14}\\[6pt]
  \rule{0.8\linewidth}{0.5pt}\\[6pt]
  {\large\bfseries Paper No. BPT-201: Thermal Physics}\\[4pt]
  {\normalsize Time: 3 Hours\hfill Maximum Marks: 70}
\end{center}
\rule{\linewidth}{0.4pt}
\smallskip
\noindent\textit{Instructions: Answer \textbf{five} questions including Question~1 which is compulsory. Symbols have their usual meanings.}
\bigskip

\noindent\textbf{Question 1.} Answer any \textbf{four} of the following: \pts{4 $\times$ 3.5 = 14}
\begin{parts}
  \item Find the probability of an oxygen molecule having its velocity in the range $99.5$ to $100.5\,\text{m/s}$ at $200\,\text{K}$. Given: molecular weight of oxygen = 32 and Avogadro number $N_A = 6.023 \times 10^{23}/\text{mol}$.
  \item Find the change in entropy of the universe in converting $1\,\text{g}$ of ice at $0^\circ\text{C}$ into steam at $100^\circ\text{C}$. (Latent heat of fusion of ice = $80\,\text{cal/g}$, latent heat of vaporization = $540\,\text{cal/g}$).
  \item Show that Gibbs function remains constant for an isobaric process carried out in such a way that the final temperature remains equal to the initial temperature at the conclusion of the thermodynamic process.
  \item Discuss whether blackbody radiation inside an enclosure remains black even after adiabatic expansion.
  \item The efficiency of a Carnot engine working in a reversible cycle between a source and a sink (temperature of sink = $127^\circ\text{C}$) is $1/6$. If the temperatures of source and sink are reduced by $100^\circ\text{C}$, find the efficiency of the Carnot engine.
\end{parts}

\medskip\rule{\linewidth}{0.2pt}

\noindent\textbf{Question 2.}
\begin{parts}
  \item How are absolute zero and unit degree of temperature defined on the thermodynamic absolute scale of temperature? \pts{5}
  \item Discuss how Andrews' experiments lead to the existence of critical constants for carbon dioxide gas undergoing liquefaction. \pts{9}
\end{parts}

\medskip\rule{\linewidth}{0.2pt}

\noindent\textbf{Question 3.}
\begin{parts}
  \item Show that the correction for volume due to the finite size of molecules for a van der Waals gas is four times the actual volume of all the molecules. \pts{5}
  \item Show that there is an increase in the entropy of the universe for an irreversible process. \pts{4}
  \item What are the Clausius and Kelvin-Planck statements of the second law of thermodynamics? Show that the Kelvin-Planck and Clausius statements of the second law of thermodynamics are equivalent. \pts{5}
\end{{parts}}

\medskip\rule{\linewidth}{0.2pt}

\noindent\textbf{Question 4.}
\begin{parts}
  \item Show that for an ideal gas, internal energy is a function of temperature only and is independent of volume, while for a van der Waals gas it depends on volume also. \pts{5}
\end{parts}
\hfill \textbf{OR}
\begin{parts}
  \item Find an expression for the fall in temperature resulting from the adiabatic stretching of a liquid film, if there is a small increase in the surface area of the film and the volume of the film remains nearly constant. \pts{4}
  \item Show that:
    \[ C_p - C_v = T \left( \frac{\partial p}{\partial T} \right)_v \left( \frac{\partial V}{\partial T} \right)_p \]
    and hence deduce that for a van der Waals gas:
    \[ C_p - C_v \approx R \left[ 1 + \frac{2a}{VRT} \right] \]
    where the symbols have their usual meanings. \pts{6}
  \item Derive the Clausius-Clapeyron equation using the appropriate Maxwell's thermodynamic relation. With the help of this relation, discuss the variation of boiling and melting points with pressure. \pts{4}
\end{parts}

\medskip\rule{\linewidth}{0.2pt}

\noindent\textbf{Question 5.}
\begin{parts}
  \item Explain the phenomenon of adiabatic demagnetization of a paramagnetic salt. How will you employ this phenomenon to produce very low temperatures? Is absolute zero attainable? \pts{8}
  \item Show that the Joule-Thomson expansion is an isenthalpic process and derive the expression for the Joule-Thomson coefficient for a van der Waals gas:
    \[ \mu_{JT} = \left( \frac{\partial T}{\partial p} \right)_H = \frac{1}{C_p} \left[ \frac{2a}{RT} - b \right] \]
    where the symbols have their usual meanings. \pts{6}
\end{parts}
\hfill \textbf{OR}
\begin{parts}
  \item Show that the perfect gas scale and the absolute thermodynamic scale are identical. \pts{6}
  \item Derive the first and second $T\,dS$ equations using Maxwell's thermodynamic relations. \pts{4}
  \item What is superfluidity with reference to liquid Helium? Discuss the phenomenon of second sound and the velocity of second sound. \pts{4}
\end{{parts}}

\medskip\rule{\linewidth}{0.2pt}

\noindent\textbf{Question 6.}
\begin{parts}
  \item Obtain the Rayleigh-Jeans formula for the distribution of intensity in the emission spectrum of a black body. Comment on its success and failure. \pts{10}
  \item Define Boyle temperature. What is the relation between Boyle temperature, critical temperature, and temperature of inversion? \pts{4}
\end{parts}

\vfill
\rule{\linewidth}{0.4pt}
\begin{center}\small
  \href{https://bhu-exam.vercel.app/}{Click here to visit our website}\\[4pt]
  \href{https://me-aryan.vercel.app/}{Click here to visit our contributor}\\[4pt]
  \href{https://quashan-me.vercel.app/}{Click here to visit our contributor}
\end{center}

\end{document}
